\section{Introduction}

Let $P$ be the self-adjoint quantum mechanical momentum operator on $L^2(\R)$, i.e., $P=-\ii\frac{\mathrm{d}}{\mathrm{d}x}$ and for $b>0$ denote by $U(b)$ the Weyl operator $U(b)=\exp(-bP)$. By using the Fourier transform
\begin{gather*}
\widehat{\psi}(k)=(\mathcal{F}\psi)(k)=\int_{\R}\e^{-2\pi\ii kx}\psi(x)\,{\rm d} x
\end{gather*}
we can write the domain of $U(b)$ as
\begin{gather*}
\dom(U(b))=\big\{\psi\in L^2(\R)\colon \e^{-2\pi b k}\widehat{\psi}(k)\in L^2(\R)\big\}.
\end{gather*}
Equivalently, $\dom(U(b))$ consists of those functions $\psi(x)$ which admit an analytic continuation to the strip $\{z = x+\ii y\in\C\colon 0<y <b\}$ such that $\psi(x +\ii y) \in L^2(\R)$ for all $0\leq y < b$ and there is a limit $\psi(x +\ii b - \ii 0) = \lim_{\eps\to 0^+}\psi (x + \ii b - \ii\eps)$ in the sense of convergence in $L^2(\R)$, which we will denote simply by $\psi(x+\ii b)$. The domain of the inverse operator $U(b)^{-1}$ can be characterised similarly.

For $b>0$ we define the operator $W_0(b)=U(b)+U(b)^{-1}=2\cosh(bP)$ on the domain
\begin{gather*}
\dom(W_0(b))=\big\{\psi\in L^2(\R)\colon 2\cosh(2\pi bk)\widehat{\psi}(k)\in L^2(\R)\big\}.
\end{gather*}
The operator $W_0(b)$ is self-adjoint and unitarily equivalent to the multiplication operator \linebreak $2\cosh(2\pi b k)$ in Fourier space. Its spectrum is thus absolutely continuous covering the interval
$[2,\infty)$ doubly.

Let $V\ge0$, $V\in L^1(\R)$ now be a real-valued potential function. The scalar inequality $2\cosh(2\pi b k)-2\ge (2\pi b k)^2$ implies the operator inequality
\begin{gather}
W_0(b)-2\ge-b^2\frac{\mathrm{d}^2}{\mathrm{d}x^2}
\label{WH}
\end{gather}
on $\dom(W_0(b))$. By Sobolev's inequality, we can conclude that the operator
\begin{gather*}%\label{WV}
W_V(b)=W_0(b)- V
\end{gather*}
is symmetric and bounded from below on the common domain of $W_0(b)$ and $V$.
We can thus consider its Friedrichs extension, which we continue to denote by $W_V(b)$. This operator acts as
\begin{gather*}
(W_V(b)\psi)(x)=\psi(x+\ii b)+\psi(x-\ii b)-V(x)\psi(x).
%\label{eq:defH}
\end{gather*}
Furthermore, by an application of Weyl's theorem (in a version for quadratic forms) and Rellich's lemma together with the fact that the form domain of $W_0(b)$ is continuously embedded in $H^1(\R)$ (as discussed at the beginning of Section~\ref{sec:sharp}) the spectrum of $W_V(b)$ consists of essential spectrum $[2,\infty)$ and discrete finite-multiplicity eigenvalues below. Details of this argument in the similar case of a Schr\"odinger operator can be found in the upcoming book \cite[Proposition 4.14]{FLW}.

We will show that the discrete spectrum satisfies a version of Lieb--Thirring inequalities for $1/2$-Riesz means.
When formulating the main result of the paper it is convenient to parametrise the eigenvalues (repeated with multiplicities) as $\lambda_j = - 2\cos(\omega_j)$, where $\omega_j\in [0,\pi]$ for
$\lambda_j\in[-2,2]$ and $\omega_j\in \ii [0,\infty)$ for $\lambda_j\le -2$.
Note that in the latter case $\lambda_j = - 2\cosh(|\omega_j|)$.

\begin{Theorem}\label{1}
Let $V\ge 0$ and let $V\in L^1(\mathbb R)$. If $W_V(b)\ge-2$, then
the discrete eigenvalues $\lambda_j=-2\cos(\omega_j)\in [-2,2)$ $($repeated with multiplicities$)$ satisfy
\begin{gather}\label{main}
\sum_{j\ge1}\frac{\sin\omega_j}{\omega_j}\le \frac{1}{2\pi b}\int_{\R}V(x) \, {\rm d} x.
\end{gather}
The constant in the inequality \eqref{main} is sharp in the sense that there is a potential
$V$ such that~\eqref{main} becomes equality.
\end{Theorem}

\begin{Remark}
Note that Theorem \ref{1} does not allow to estimate eigenvalues below $-2$. In~fact, from the proof of this theorem, the case of one eigenvalue below $-2$ could be included in the inequality \eqref{main}. We expect that the inequality holds true for all eigenvalues below $-2$. However, the method we use in the proof prevents us from including all eigenvalues due to oscillating properties of the resolvent $(W_0(b) - \lambda)^{-1}$ for $\lambda<-2$.
\end{Remark}

Lieb--Thirring inequalities were first established for Schr\"odinger operators in \cite{LT}. For a~one-dimensional Schr\"odinger operator ${-}\frac{\mathrm{d}}{\mathrm{d}x^2}-V$ on $L^2(\R)$ with negative eigenvalues \mbox{$\mu_1\le\mu_2\allowbreak\le\!\cdots\!<0,\!$} these bounds state that for any $\gamma\ge1/2$ there is a constant $L_\gamma>0$ such that
\begin{gather}
\sum_{j\ge 1}|\mu_j|^{\gamma}\le L_{\gamma}\int_{\R}V(x)^{\gamma+1/2}\, {\rm d} x
\label{LTSch}
\end{gather}
for all $V\ge0$, $V\in L^{\gamma+1/2}(\R)$. The condition $\gamma\ge1/2$ is optimal. Inequality \eqref{WH} implies that
\begin{gather}
\sum_{j\ge 1}|\lambda_j-2|^{\gamma}\le \frac{L_{\gamma}}{b}\int_{\R}V(x)^{\gamma+1/2}\, {\rm d} x
\label{LTS}
\end{gather}
for all eigenvalues $\lambda_j\le 2$ of $W_V(b)$. Under the additional assumption $W_V(b)\ge-2$, our bound~\eqref{main} presents an improvement of \eqref{LTS} for $\gamma=1/2$. This can be seen from the fact that for $\gamma=1/2$ the sharp constant in \eqref{LTSch} is given by $L_{1/2}=1/2$ \cite{HLT} and from the strict inequality
\begin{gather*}
|\lambda_j-2|^\frac12=|2\cos\omega_j+2|^\frac12<\frac{\pi\sin\omega_j}{\omega_j}
\end{gather*}
for $\omega_j\in[0,\pi)$.
The difference of the terms above vanishes as $\omega_j\to\pi$, implying that \eqref{LTS} is asymptotically optimal for small coupling. While the necessity of $\gamma\ge1/2$ in the Lieb--Thirring inequality for Schr\"odinger operators does not allow us to conclude that \eqref{LTS} fails for $0\le\gamma<1/2$, we will prove the following.

\begin{Theorem}\label{th:res}
Let $b>0$. If $V\in L^1(\R)$ with $\int_{\R}V\, {\rm d} x>0$, then $W_V(b)$ has at least one eigenvalue below $2$. Furthermore, if $0\le\gamma<1/2$, then there is no constant $L_\gamma$ such that \eqref{LTS} holds for all compactly supported $V$. This conclusion holds even under the assumption that $W_V(b)\ge-2$.
\end{Theorem}
